\section{Introducción}

Las metaheurísticas representan una de las áreas más interesantes de la inteligencia artificial. 
Su importancia radica en la capacidad de ofrecer estrategias para hallar buenas soluciones en 
espacios de búsqueda complejos, donde algoritmos clásicos resultarían computacionalmente inviables. 

En el presente trabajo se analiza la eficacia y el comportamiento de tres algoritmos que aplican distintos 
enfoques, los dos primeros inspirados en la naturaleza y el último en la lógica matemática, para explorar 
estos espacios de búsqueda:

\begin{itemize}
    \item \textbf{Optimización Basada en Biogeografía (BBO):} Este algoritmo se inspira en el estudio de la 
    distribución geográfica de las especies biológicas. 
    Su mecanismo principal se basa en el intercambio de información mediante los operadores de migración
    y mutación, permitiendo que las soluciones con mayor índice de idoneidad (HSI) compartan sus atributos 
    con las de menor calidad para mejorar la población global.
    
    \item \textbf{Sistema de Hormigas Elitista (EAS):} Basado en el comportamiento de las colonias de hormigas, 
    este método utiliza el concepto del rastro de feromonas para establecer rutas eficientes. 
    Mediante un proceso de retroalimentación positiva y una estrategia elitista, el algoritmo refuerza los caminos 
    que componen las mejores soluciones encontradas, facilitando la convergencia en problemas complejos como el de 
    asignación generalizado.
    
    \item \textbf{Fuzzy Adaptive Neighborhood Search (FANS):} Se presenta como un entorno de búsqueda adaptativo 
    que integra la lógica difusa para cualificar las soluciones. 
    Su gran versatilidad se encuentra en el uso de una valoración difusa ($\mu$) que permite modular el comportamiento 
    del algoritmo, permitiéndole emular diversas metaheurísticas clásicas o adaptarse a perfiles de decisión específicos, 
    como criterios conservadores o de diversificación.

\end{itemize}

A través de la implementación de estos métodos, se pretende comprender no solo su fundamentación teórica, sino también su rendimiento práctico ante problemas \textit{benchmark} y escenarios de optimización multiobjetivo.